% ============================================================
% Section unique des limites (agent 7)
% À insérer dans chapters/chapitre6_approfondissements_scientifiques.tex
% (chapitre 5 « Validité scientifique » du PDF), après la section
% « Conclusions soutenues » (fin de fichier).
% Chaque valeur est tracée dans PATCH_NOTES.md, partie B, tableau B.0.
% ============================================================

\section{Limites du travail}
\label{sec:limites}

Cette section rassemble les limites du travail. Les autres chapitres y
renvoient sans les reprendre.

\paragraph{Données et généralisation.}
CICIDS2017 et CSE-CIC-IDS2018 reposent sur des scénarios construits et sur
des caractéristiques de flux déjà extraites. Sur CICIDS2017, le F1 du Random
Forest passe de $0{,}995142$ en séparation aléatoire à $0{,}157744$ en
holdout par scénario : le transfert vers une attaque absente de
l'entraînement n'est pas acquis. Sur CSE-CIC-IDS2018, \texttt{Dst Port} seul
atteint un F1 de $0{,}999720$ et le retrait des dix premières
caractéristiques ramène le F1 à $0{,}657859$ ; un artefact lié aux journées
retenues ou au processus de génération du jeu ne peut être exclu. Aucun de
ces résultats ne se transpose sans nouvelle évaluation à d'autres capteurs,
topologies ou périodes d'observation.

\paragraph{Taille des échantillons.}
Le test HDFS compte 29 blocs positifs sur $2\,000$ ; l'intervalle bootstrap
traduit cette faiblesse sans corriger un éventuel biais de sélection du
corpus. Sur BGL, le groupe des templates connus ($2\,029$ unités) ne contient
aucune anomalie : le rappel des anomalies à template connu ne peut pas être
estimé, et $89{,}8081\,\%$ des unités de test relèvent d'un template inconnu
de l'entraînement. L'évaluation open-set du routeur repose sur trois fichiers
inconnus et 28 fichiers connus, et la \texttt{decision\_margin} n'est pas une
probabilité calibrée. La voie Apache repose sur une seule fixture.

\paragraph{Grandeurs non mesurées.}
Les sources de la campagne multi-source ne partagent pas de vérité terrain
homogène : aucune exactitude globale n'est calculée et le F1 obtenu sur
Linux/auth reste propre à cette source. La couverture multi-format mesure la
lecture, le parsing et la normalisation des unités présentes, non la
complétude des champs ni la conservation intégrale du contenu brut. La
corrélation est évaluée sur un corpus synthétique. L'utilisabilité du tableau
de bord, le temps d'investigation et la charge cognitive n'ont fait l'objet
d'aucune étude auprès d'analystes.

\paragraph{Distribution et exploitation.}
La validation distribuée porte sur deux machines virtuelles et un Redis
central pendant $3\,605{,}110$~s. Ni panne du broker, ni partition réseau, ni
fonctionnement multi-site, ni charge industrielle, ni vieillissement sur
plusieurs jours ne sont testés. Les mesures CPU et mémoire décrivent des
processus locaux et non une infrastructure complète. TLS, contrôle d'accès,
gestion des secrets, journalisation protégée contre l'altération, réplication,
sauvegarde et observabilité distribuée ne sont pas évalués.

\paragraph{Mécanismes adaptatifs.}
Les poids de la fonction d'utilité sont fixés et non appris
(section~\ref{sec:spec-utilite}). La mémoire modifie certaines décisions, mais
le benchmark A--D ne montre pas de gain systématique de débit, de latence ou
d'équité ; ce benchmark porte sur des micro-tâches, et l'effet de tâches plus
coûteuses reste une hypothèse. La poursuite d'objectifs, l'enchaînement, la
délégation et la reprise sont vérifiés par des tests fonctionnels, sans
campagne quantitative mesurant leur coût ou leur bénéfice. La mise à jour
contrôlée des modèles dispose d'une implémentation et d'un mode
\textit{dry-run}, mais d'aucun cycle complet et répété de promotion et de
rejet suivi d'une mesure indépendante (tableau~\ref{tab:statut-model-update}).

\paragraph{Statut du système.}
Ariel \logminer{} est en conséquence un prototype de laboratoire. Il ne
constitue ni une plateforme opérationnelle de SOC ni un système autonome prêt
pour la production.
